\documentclass[languages=english,francais]{resume_cover_letter}

% infos
\firstname{Gaël}
\familyname{Foppolo}  
\title{Étudiant ingénieur @ Polytech Marseille}
\address{}{\faHome~Marseille -- France}   
\email{me@gaelfoppolo.com}                      
\homepage{www.gaelfoppolo.com}
\mobile{+33 6 28 91 02 66}
\social[github]{gaelfoppolo}
\social[linkedin]{gaelfoppolo} 
\social[twitter]{gaelugio} 
\extrainfo{Permis B}
\photo[95pt][0pt]{img/photo}	
\quote{À la recherche d'opportunités de stage de fin d'études (6 mois) dans la Data Science, le DevOps ou iOS}

\begin{document}

\maketitle

\vspace*{-7mm}

\section{\faBriefcase~Expériences}

%\vspace*{-1mm}

\tlcventry{2017}{2018}{Self-service data analytics}{CapGemini}{Marseille}{\url{www.capgemini.com}}{\textit{Technologies : Hadoop, Spark, AWS}}

\vspace*{2mm}

\tldatecventry{2017}{Stage - Ingénieur iOS \faApple~(3 mois)}{Care Labs}{Montpellier}{\url{www.carelabs.co}}{L'application \textit{CareBook} est un carnet de santé personnel, permettant de consulter ses rendez-vous, stocker ses documents ou mettre à jour son dossier médical. Ma mission a été la création complète de l'application : conception, développement, tests, recettage et déploiement.
}

\vspace*{2mm}

\tldatecventry{2016}{Stage - Développeur iOS \faApple~(2 mois)}{Lundi Matin}{Montpellier}{\url{www.lundimatin.fr}}{\textit{RoverCash} remplace la caisse enregistreuse classique, transformant tablette et smartphone en véritable outil de gestion. L'application offre, en plus des fonctionnalités traditionnelles : annuaire, catalogue, vente, paiement, statistiques, etc.\\
	Mon objectif principal durant ce stage a été l’internationalisation (i18n) et la localisation
(l10n) de l’application.
\newline{\url{https://itunes.apple.com/app/id1061525820}}}

\vspace*{2mm}

\tldatecventry{2015}{Stage - Développeur iOS \faApple~(3 mois)}{KeepCore}{Montpellier}{\url{www.keepcore.com}}{L'application \textit{HandiCarParking} a pour objectif de localiser rapidement les places de stationnement réservées aux personnes handicapées dans le monde entier. Elle se base sur les données d'\textit{OpenStreetMap}.\\
	Ma mission a été la création complète de l'application, des spécifications au déploiement en passant par le développement.
\newline{\url{https://github.com/gaelfoppolo/handicarparking}}}

%\vspace*{-2mm}

\section{\faGavel~Projets}

%\vspace*{-1mm}

\tldatecventry{2017}{AnalytiCeM}{Polytech}{Marseille}{\url{github.com/gaelfoppolo/AnalytiCeM}}{
AnalytiCeM a pour objectif de récupérer et d'analyser des données cérébrales, provenant d'un casque Muse. Pour cela, on utilise la notion de « session », séance pendant laquelle on récupère des données contextuelles (comme la météo) que l’on enrichit avec les signaux récupérés par le casque. On peut ainsi inférer des comportements en fonction du contexte.\newline{\url{} \hfill \textit{Technologies : Swift \& Realm}}}

\vspace*{2mm}

\tldatecventry{2016}{Locomotor}{Polytech}{Marseille}{\url{github.com/gaelfoppolo/locomotor}}{
Locomotor est un comparateur de moyens de transport réels et fictifs. Il se base sur des critères prédéfinis que l'utilisateur peut sélectionner. Il génère une liste de véhicules qui correspondent \og le mieux \fg à la demande de l'utilisateur.
\newline{\url{} \hfill \textit{Technologies : Java \& MongoDB}}}

\vspace*{2mm}

\tldatecventry{2015}{Apprentissage supervisé}{Polytech}{Marseille}{\url{github.com/gaelfoppolo/examples-learning}}{Ce projet a pour but de trouver les similarités dans un jeu de données formatées et de proposer une ou plusieurs descriptions caractéristiques qui apportent le plus d'informations. L'application produit des solutions de la forme : \textit{\og un carré rouge de taille 8 à 10 est sur un triangle bleu ou vert de taille 5, lui-même à côté d'un cercle rouge de taille 2 à 4\fg}
\newline{\url{} \hfill \textit{Technologies : C \& C++}}}

%\vspace*{-3mm}

\section{\faGraduationCap~Formations}

%\vspace*{-1mm}

\tlcventry{2015}{2018}
{Ingénieur informatique}{Polytech}{Marseille}{}{}

\tlcventry{2013}{2015}
{DUT Informatique}{Institut Universitaire de Technologie}{Montpellier}{}{}

\vspace*{-3mm}

\section{\faNewspaperO~Publication}

\vspace*{-1mm}

\begin{otherlanguage}{english}

\tldatecventry{2015}{Contextual Sequential Pattern Mining in Games: Rock, Paper, Scissors, Lizard, Spock}{Dumartinet, J., Foppolo, G., Forthoffer, L., Marais, P., Croitoru, M., \& Rabatel, J.}{Research and Development in Intelligent Systems XXXII (pp. 375-380)}{Springer International Publishing}{\url{https://doi.org/10.1007/978-3-319-25032-8_29}}

\end{otherlanguage}

\vspace*{-3mm}

\section{\faLanguage~Langues}

\vspace*{-1mm}

\begin{minipage}{\linewidth}
  \begin{multicols}{2}
    \centering{\textbf{Français} -- \faStar~\faStar~\faStar~\faStar~\faStar\\}
    \vspace*{1mm}
    \centering{\textit{Projet Voltaire -- 80\%}}\\
    \columnbreak
    \centering{\textbf{Anglais} -- \faStar~\faStar~\faStar~\faStar~\faStarO\\}
    \vspace*{1mm}
    \centering{\textit{TOEIC -- 935}}\\
  \end{multicols}
\end{minipage}

%\vspace*{-2mm}

\section{\faGears~Compétences}

%\vspace*{-2mm}

\begin{minipage}{\linewidth}
  \begin{multicols}{4}
  \centering{\textbf{Environnements\\}}
    \vspace*{2mm}
        
        	%\vspace*{3mm}
        \centering{macOS~~\textsc{\faStar~\faStar~\faStar~\faStarHalfO~\faStarO}\\}
        \centering{iOS~~\textsc{\faStar~\faStar~\faStar~\faStarHalfO~\faStarO}\\}
    	\centering{UNIX~\textsc{\faStar~\faStar~\faStarHalfO~\faStarO~\faStarO}\\}
    	%\centering{MongoDB~\textsc{\faStar~\faStar~\faStarO~\faStarO~\faStarO}\\}
    	    	
    \columnbreak
    \centering{\textbf{Langages\\}}
    \vspace*{2mm}
   
   \centering{Swift~~\textsc{\faStar~\faStar~\faStar~\faStarO~\faStarO}\\}    
    \centering{Java~~\textsc{\faStar~\faStar~\faStar~\faStarO~\faStarO}\\}    
    \centering{SQL~~\textsc{\faStar~\faStar~\faStarHalfO~\faStarO~\faStarO}\\}
    	
       \columnbreak
    \centering{\textbf{Outils\\}}
    \vspace*{2mm}
    
    	%\vspace*{3mm}
    	
    	\centering{Git~~\textsc{\faStar~\faStar~\faStar~\faStarHalfO~\faStarO}\\}
    	\centering{Slack/Trello~\textsc{\faStar~\faStar~\faStar~\faStar~\faStarHalfO}\\}
    	\centering{UML~~\textsc{\faStar~\faStar~\faStar~\faStarO~\faStarO}\\}
    	
    	\columnbreak
    \centering{\textbf{Personnelles\\}}
    \vspace*{2mm}
        
        	%\vspace*{3mm}
        \centering{Leadership\\}
        \centering{Réactivité\\}
    	\centering{Polyvalence\\}
    	%\centering{MongoDB~\textsc{\faStar~\faStar~\faStarO~\faStarO~\faStarO}\\}

    
  \end{multicols}
\end{minipage}

%\vspace*{-1mm}

\section{\faBeer~Autres}

%\vspace*{-1mm}

\tikzset{
    cercle/.pic={
      \node [draw, thick, circle, minimum width=10pt] {\tikzpictext};
    },
  }
\hspace*{1mm}
\begin{minipage}{\linewidth}
  \begin{tikzpicture}
  	\pic [pic text={\LARGE \faCoffee}] {cercle};
    \node[draw=none] at (0,-0.9) {Thé};
    \pic [pic text={\LARGE \faTelevision}] at (20mm,0) {cercle};
    \node[draw=none] at (2,-0.9) {Séries};
    \pic [pic text={\LARGE \faSpotify}] at (40mm,0) {cercle};
    \node[draw=none] at (4,-0.9) {Spotify};
    \pic [pic text={\LARGE \faLeaf}] at (60mm,0) {cercle};
    \node[draw=none] at (6,-0.9) {Jardinage};
    \pic [pic text={\LARGE \faApple}] at (80mm,0) {cercle};
    \node[draw=none] at (8,-0.9) {Apple};
    \pic [pic text={\includegraphics[width=.03\textwidth]{img/running.png}}] at (100mm,0) {cercle};
    \node[draw=none] at (10,-0.9) {Running};
    \pic [pic text={\LARGE \faFirefox}] at (120mm,0) {cercle};
    \node[draw=none] at (12,-0.9) {Firefox};
    \pic [pic text={\LARGE \faMoonO}] at (140mm,0) {cercle};
    \node[draw=none] at (14,-0.9) {Nuit};
    \pic [pic text={\LARGE \faGithub}] at (160mm,0) {cercle};
    \node[draw=none] at (16,-0.9) {Open source};
  \end{tikzpicture}
\end{minipage}

\end{document}